% ============================================================================
% Introduction (from Arce's gBKMR_Nature.docx)
% ============================================================================

\section*{Introduction}

Exposure to complex mixtures of environmental chemicals and other toxic substances is a pervasive issue of public health. Many of these substances have been linked to adverse health effects, including toxic metals, air pollution, PFAS, and, more recently, micro and nanoplastics. The quantification of the joint health effects of these substances over an individual lifetime is of interest for prevention and policy making, however, it poses great statistical challenges.

Approaches for the analysis of multiple exposures are now available. Billionnet et al.\cite{billionnet2012} provided a review of existing methods to analyze the health effects of exposure to multipollutant mixtures. Depending on the question and the setting, Weighed Quantile Sum Regression (WQSR),\cite{carrico_wqs_2015} Bayesian Kernel Machine Regression (BKMR),\cite{bobb_bkmr_2015} and the Generalized Additive Model (GAM) have been shown to work well in practice. These approaches have even been extended to the Bayesian setting, including Bayesian Weighted Sums\cite{colicino_bwqs_2020} and Bayesian Generalized Additive models. However, these approaches are limited to analyzing data at a single time point, while longitudinal studies that collect data sequentially at multiple time points are common in fields such as epidemiology, sociology, and econometrics, and provide a more realistic picture of the long-term health effects of the mixture. Alternative models that characterize change for longitudinal data, such as Generalized Estimating Equations (GEE) or Generalized Linear Mixed Models (GLMM), do not afford a causal interpretation, unless appropriate strategies to accommodate time-dependent confounding are adopted.

Causal inference methods center on the premise that data from an observational study can be interpreted as originating from a randomized clinical trial in which the treatment was randomly assigned based on (conditional on) measured covariates. Time-varying confounders are factors that have the potential to confound the causal relationship between a time-varying exposure variable and outcome and that themselves change over time and are evaluated frequently throughout the investigation. When confounders both influence and are influenced by time-varying exposures, standard regression adjustment leads to bias. Approaches to account for time-varying confounding are required, such as g-methods (e.g., standardization or weighting).

In environmental health sciences, there is a growing demand for approaches that can handle longitudinal data in the presence of multiple correlated exposures. However, there are methodological challenges that lead to uncertainty regarding how to translate the results of environmental mixture studies into policy recommendations. For example, incorrect adjustment for confounding can lead to bias, and most approaches focus on flexible modeling of environmental mixture effects, leaving confounders adjusted assuming a linear relationship. This is problematic because confounders may interact with the environmental mixture, leading to biased results. Approaches that allow for lagged exposures have been developed under the BKMR framework, however, none of them accommodates time-dependent confounding.

Quantifying the joint effect of environmental mixtures over time is crucial to determine intervention timing. However, currently, there is a lack of statistical approaches that can simultaneously address time-varying confounding, flexible modeling, and variable selection when assessing the effect of multiple, correlated, and time-varying exposures. To bridge this gap, we have developed a causal inference method, g-BKMR, which allows for the estimation of nonlinear, non-additive effects of both time-varying exposures and confounders while also enabling variable selection. By combining the popular BKMR method and the g-formula, g-BKMR offers a unique approach to examine the non-linear and non-additive effects of exposure mixtures when their effects vary over time. We present a simulation study and an application of our method to evaluate the association between a mixture of 6 urinary metals (arsenic, cadmium, molybdenum, selenium, tungsten, and zinc) and fasting blood glucose in the Strong Heart Study, a prospective cohort of American Indians with high rates of cardiovascular disease and low to moderate exposure to metals.
